% Generated from paper/full_draft. Edit the Markdown source or migration script.

\section{Conclusion}
\label{sec:09-conclusion:conclusion}

This paper presents IntentRoute, a feedback-adaptive confidence-gated route and calibrated-budget controller motivated by local relevance structure in vertical-domain data. In the evaluated retrieval-backed QA implementation, geometry defines reproducible cluster-local routes, trust-weighted LinUCB updates route confidence, dense and BM25 provide rescue paths, and a calibrated policy separately controls the final evidence-context budget.

The evaluation spans seven dataset settings with different evidentiary roles. The main full-stack evidence comes from LoTTE technology/search at 100k to 638k corpus chunks. Under calibration/test budget selection, calibration-eligible operating points at 100k, 200k, and 638k reduce final LLM evidence-context input tokens by 6-18\%; the original 400k point remains calibration-ineligible. A normalized five-fold follow-up at 400k yields 14.50\% mean saving with no mean Hit change, while retaining policy instability and no strict seed-level non-inferiority. Across these scales, IntentRoute avoids the larger $\mathrm{Hit@10}$ losses of dense-only adaptive truncation, while strict seed-level non-inferiority remains scale-dependent. A conservative confidence-only policy remains a stable 4.7-5.3\% saving baseline. Split sensitivity checks strengthen the 200k and 638k operating points while showing that 100k and especially 400k policy selection is more partition-dependent. Additional diagnostics and controls show that local geometry provides useful route signal over random routing and trust-weighted feedback improves route confidence over no-feedback controls, without implying that either alone explains fused quality. Matched BGE/E5 comparisons retain near-dense retrieval quality with about 12\% context saving, while a BGE quality-first point demonstrates frontier tunability. A 300-query evaluation with three LLM judges finds approximately 6-12\% matched context savings without a statistically detectable correctness change, while exposing method-dependent faithfulness effects and retaining the lack of strict answer-level non-inferiority. LoTTE science/search provides cross-domain ranking support with a clear compression-calibration boundary. Hard-case recovery experiments further show that simulated feedback can repair part of the tail failures caused by aggressive context compression. PubMedQA and CovidQA-RAG extend the evidence-retrieval transfer checks to biomedical QA under near-ceiling and more discriminative dense baselines, respectively; Banking77 extends the feedback-adaptation check to banking-intent routing; and eManual and CUAD expose duplicate-text and sparse-ground-truth limits. These supporting settings broaden the transfer, mechanism, and boundary evidence without being treated as equivalent replications of the LoTTE quality-efficiency frontier.

Strong post-retrieval baselines refine rather than weaken the conclusion. Sentence-level MMR and Selective Context-lite are effective shared downstream compressors, and cross-encoder reranking can improve top-ranked evidence support. However, reranking alone can increase final context tokens, and same-budget reranking does not uniformly dominate the calibrated IntentRoute policies. These results support a layered interpretation: candidate generation, reranking, compression, route control, and budget calibration are separate system functions that can be composed.

The result is intentionally bounded. IntentRoute is not a universal dense replacement, a universal compressor replacement, or a universal reranker replacement, and it does not prove that geometry alone solves retrieval. Dense retrieval remains an important recall floor. The contribution is a calibrated controller that combines geometry- and feedback-informed route control with a separately calibrated final context budget, trading compact context against retrieval risk while remaining compatible with late reranking and prompt compression. The manifold hypothesis remains the motivation for local route structure, not a theorem-level claim.
